\subsection{동결된 \texttt{reveal\_hint} 계약과 공개 증거}
이 산출물은 동결 기대값 패킷 \texttt{sl-npc-disclosure-expectations-v1} (v1.0.0)과 네 개 실제 결과 파일만 읽는다. 패킷 SHA-256 \texttt{9920aa834910a3ee\allowbreak{}35ed4446de4fdba4\allowbreak{}2bb67dda71b07a93\allowbreak{}bccfc7a0b2e09028}는 sidecar와 대조했다. 선언식 \texttt{1 eligible actor x 4 declared fact IDs x 3 declared stage values x 2 knowledge values x 2 prior-disclosed values}의 정상 분모는 48개이며, 경계 4개와 잘못된 입력 8개는 제외한다. Trace 기대값은 gate 출력이 아닌 packet case ID로 대조한다.

내용 우선 설명은 typed input(연산·행위자·fact·stage·knowledge·prior-disclosed)을 ontology entities/relations와 authored rules에 통과시켜 check, allowed/denied 결정, canonical state effect를 기록한다. 허용된 최초 표시는 요청 fact만(조석 단서의 선언된 relationship-memory 항목도) 추가하고 stage를 보존하며 revision을 한 번 증가시킨다. 거부·기공개 재표시는 canonical state를 보존한다. 설명용일 뿐 hidden fact·자연어 함의·live capture를 주장하지 않는다.

온톨로지 선언은 \texttt{sl-npc-disclosur\allowbreak{}e-expectations-v\allowbreak{}1:runtime-disclo\allowbreak{}sure} / v1.0.0, provenance는 \texttt{frozen\_expectation\_packet}이다. 범위는 네 canonical fact ID(dock\_saved, keeper\_betrayal, lighthouse\_dark, tide\_marks\_hint)와 행위자 \texttt{captain\_mira}뿐이다. 규칙 ID: DISC-001, DISC-002, DISC-003, DISC-004, DISC-005, DISC-006, DISC-007; 승인 template은 \texttt{mira-safe-deflection-v1} 및 \texttt{sl-rc-003-authorized-hint}뿐이다.

Python evaluator는 같은 사례·기대 객체를 세 격리 arm에 materialize하고 semantic outcome, check trace, error, 전후 state SHA-256, template ID, replay equality를 기록한다. Godot은 별도 game-execution 경계에서 기존 \texttt{reveal\_hint} writer와 canonical mutation을 소유하며, 이 producer는 state를 재구성하거나 대체하지 않는다. 연결은 operation \texttt{reveal\_hint}, scenario \texttt{sealed-lighthouse-v1}, actor \texttt{captain\_mira}, bound source hash로 추적한다.
실행 시점 관측 hash는 현재 source와 일치한다. 동결 패킷 baseline과 다른 source는 3개(game-track/godot/scripts/game3d/game\_3d.gd, game-track/godot/scripts/sealed\_lighthouse\_machine.gd, llm-wiki/wiki/reports/2026-10-22-ontology-signals-interview.md)이며, 패킷 hash는 수정하지 않고 이 차이를 공개한다.

표~\ref{tab:disclosure-arms}는 같은 60개 사례의 세 arm을 정상/경계/잘못된 입력 분모로 보고하며, unchecked 위반은 의도된 격리 관측이다. 검증된 general-rule과 ontology-gate arm은 저작된 사례 전체에서 결과가 동일하다(60/60). rule-change fixture 집합에는 두 arm을 가르는 discriminating case가 없으며, 2개 fixture의 4개 paired case가 모두 일치한다. 이는 공유 동결 규칙 아래 authorization-conformance limit이며 일반적 ontology equivalence 또는 ontology superiority의 증거가 아니다. 그림~\ref{fig:disclosure-signals}는 E/R(기대 일치·재생 사례), C(check-record PASS/FAIL), NR(보고된 check가 없는 사례)를 구별한다. NR은 skipped/pass/out-of-scope가 아니다. 표~\ref{tab:disclosure-fixtures}는 rule-change fixture와 paired case를 disposable clone에 매핑하며 shipped scenario는 바꾸지 않는다.

주장은 bounded engineering conformance에 한정된다. human gold, access-separated holdout, efficacy, superiority, generalization, player benefit, model latent state, cost는 주장하지 않는다. 정확한 입력 경로·hash·재생성 명령은 \path{paper/latex/generated/disclosure_producer_manifest.json}에 결속된다.
